\documentclass[conference]{IEEEtran}
% ---------------------------------------------------------------------------------------------
% IEEE conference version (target: 9 pages including references).
% Experimental values come from numbers.tex (written by src/export_latex.py from
% results/metrics.json); missing values print as a red [TBD].
% ---------------------------------------------------------------------------------------------
\IEEEoverridecommandlockouts
\usepackage{cite}
\usepackage{amsmath,amssymb}
\usepackage{graphicx}
\usepackage{booktabs}
\usepackage{array}
\usepackage[table]{xcolor}
\usepackage{enumitem}
\usepackage{textcomp}
\usepackage[hidelinks]{hyperref}

% ----------------------------------------------------------------- result macros
% Real results (numbers.tex) take precedence; illustrative dummies are used only if it is absent.
\newcommand{\dummyval}[1]{\textcolor{blue!70!black}{#1}}% set to {#1} to print dummies in black
\InputIfFileExists{numbers.tex}{}{\InputIfFileExists{dummy_numbers.tex}{}{}\InputIfFileExists{partial_real_numbers.tex}{}{}}% real values available so far override the dummies
\newcommand{\tbd}{\textcolor{red}{\textbf{[TBD]}}}
\makeatletter
\newcommand{\res}[2]{\@ifundefined{res@#1@#2}{\tbd}{\@nameuse{res@#1@#2}}}
\newcommand{\respm}[2]{\@ifundefined{res@#1@#2}{\tbd}{\@nameuse{res@#1@#2}%
  \@ifundefined{resstd@#1@#2}{}{\,{\tiny$\pm$\@nameuse{resstd@#1@#2}}}}}
\newcommand{\teach}[1]{\@ifundefined{teach@#1}{\tbd}{\@nameuse{teach@#1}}}
\newcommand{\envv}[1]{\@ifundefined{env@#1}{\tbd}{\@nameuse{env@#1}}}
\makeatother
\newcommand{\interp}[1]{\textcolor{red}{\textit{[Pending results: #1]}}}
\newcommand{\figorbox}[2]{\IfFileExists{figures/#1}{\includegraphics[width=#2]{figures/#1}}%
  {\fbox{\parbox[c][3cm][c]{\dimexpr#2-2\fboxsep-2\fboxrule\relax}{\centering\textcolor{red}{\textbf{[Figure pending]}}\\\scriptsize\texttt{\detokenize{#1}}}}}}
\newcommand{\Lq}{\mathcal{L}}
% ----------------------------------------------------------------- display names (rename here)
\usepackage{xspace}
\newcommand{\method}{\textup{\textsc{NumBridge}}\xspace}
\newcommand{\vBase}{\textup{\textsc{Base}}\xspace}                      % original all-MiniLM-L6-v2
\newcommand{\vParse}{\textup{\textsc{Base+Parse}}\xspace}                % original + symbolic parser channel
\newcommand{\vParseOnly}{\textup{\textsc{ParseOnly}}\xspace}             % parser alone (reference)
\newcommand{\vExact}{\textup{\textsc{NumBridge-Exact}}\xspace}           % exact (mathematical) supervision
\newcommand{\vMatched}{\textup{\textsc{NumBridge-Matched}}\xspace}       % exact labels on the reasoner's pair budget
\newcommand{\vRaw}{\textup{\textsc{NumBridge-Raw}}\xspace}               % raw LLM-reasoner labels
\newcommand{\vVerified}{\textup{\textsc{NumBridge-Verified}}\xspace}     % LLM-reasoner labels, verified
\newcommand{\vExactParse}{\textup{\textsc{NumBridge-Exact+Parse}}\xspace}% exact + parser channel
\newcommand{\vDual}{\textup{\textsc{NumBridge-Dual}}\xspace}             % learned numerical subspace, no parser
\newcommand{\vAbs}{\textup{\textsc{NumBridge-Abs}}\xspace}               % exact, absolute-distance target


\newcommand{\Ea}{exp1_original_minilm}
\newcommand{\Eb}{exp2_parser_plus_original}
\newcommand{\Ep}{ref_parser_only}
\newcommand{\Ec}{exp3_math_full}
\newcommand{\Ecb}{exp3b_math_budget}
\newcommand{\Eda}{exp4a_teacher_raw}
\newcommand{\Edb}{exp4b_teacher_verified}
\newcommand{\Ee}{exp5_hybrid_finetuned_plus_parser}
\newcommand{\Ef}{exp6_learned_projection}
\newcommand{\Eabs}{exp3_math_full_absdist}

\begin{document}

\title{\textsc{NumBridge}: Can Small LLMs Teach\\Sentence Encoders to Understand Numbers?}

\author{\IEEEauthorblockN{Tariq Mahmood}
\IEEEauthorblockA{\textit{AI \& Data Science Lab, DIMES} \\
\textit{University of Calabria}\\
Rende (CS), Italy}}

\maketitle

\begin{abstract}
Dense sentence encoders underpin retrieval-augmented generation, semantic search and entity resolution, yet they
largely treat numbers as surface tokens: \emph{10\,kg} and \emph{10{,}000\,g} may be far apart, while \emph{10\,kg} and
\emph{10\,m} look alike. We ask whether the numerical competence of a small instruction-tuned LLM can be
\emph{transferred} into a lightweight encoder (all-MiniLM-L6-v2) by knowledge distillation, and how one would know
that the encoder has learned numbers rather than benefiting from a symbolic parser. We contribute (i) a formulation of
numerical desiderata for cosine-similarity spaces, including what symmetric similarity cannot express; (ii) a
distillation protocol with a reasoner (Qwen2.5) verified against deterministic ground truth, paired with a
budget-matched reasoner-free control that isolates the reasoner's contribution; (iii) attribution controls that separate
embedding-internal learning from parser-driven gains, namely a parser-free numbers-in-words slice and a learned
numerical projection head; and (iv) a leakage-controlled benchmark of 5{,}010 statements over six physical dimensions,
with affine temperature handling. We pre-declare hypotheses and decision rules and report preliminary results.
\end{abstract}

\begin{IEEEkeywords}
sentence embeddings, numeracy, knowledge distillation, large language models, semantic retrieval
\end{IEEEkeywords}

\ifdefined\ResultsAreDummy
\noindent\fcolorbox{orange!80!black}{orange!10}{\parbox{\dimexpr\columnwidth-2\fboxsep-2\fboxrule}{\footnotesize\textbf{Illustrative dummy values:}
numbers in \dummyval{blue} are layout placeholders, \emph{not results}. Numbers in black (reasoner accuracy,
hardware) are real, from the run in progress. All are replaced automatically by \texttt{numbers.tex}.}}
\fi
\ifdefined\ResultsAreMock
\noindent\fcolorbox{red}{red!8}{\parbox{\dimexpr\columnwidth-2\fboxsep-2\fboxrule}{\footnotesize\textbf{Pipeline-check build:}
numbers come from a random tiny encoder and a mock reasoner; they are not findings.}}
\fi

% =============================================================================================
\section{Introduction}
% =============================================================================================
Take ``The package weighs 10\,kg'', ``The package weighs 10{,}000\,g'' and ``The package is 10\,m long''. The first two
state the same quantity; the third a different dimension. A general-purpose sentence encoder~\cite{reimers2019sbert}
typically scores the third at least as close to the first as the second: the wording overlaps, the digits match, and
the encoder has no notion of unit conversion. Probing studies find magnitude only partly and brittlely decodable from
embeddings~\cite{naik2019numeracy,wallace2019numbers,zhang2020scales,thawani2021survey}. In retrieval-augmented
generation~\cite{lewis2020rag} and in entity resolution over numeric attributes, such errors propagate silently.

Small instruction-tuned LLMs such as Qwen2.5~\cite{qwen2025qwen25} often answer ``is 10\,kg the same as 10{,}000\,g?''
correctly, but are too costly to run inside every similarity computation. \textbf{Numerical Intelligence Transfer}
distils this behaviour into a compact encoder~\cite{hinton2015distilling,reimers2020multilingual} so that it
is paid for once, at training time. We call the resulting framework \method: it bridges the numerical competence of an LLM \emph{reasoner} and the efficiency of a compact retrieval \emph{encoder}.

The idea is simple; evaluating it honestly is not. Three confounds can make a transfer method look successful when it
is not. \textbf{(C1) Cosine is not numeracy}: higher similarity on equivalent pairs is meaningless if all similarities
rise. \textbf{(C2) The parser confound}: any pipeline that extracts numbers symbolically excels on parsable text,
whether or not the embedding learns anything. \textbf{(C3) The verification paradox}: if reasoner labels are filtered
against ground truth---good practice---the surviving labels \emph{equal} that ground truth on every determinable
relation, so what does the reasoner add?

\textbf{Contributions.}
\begin{itemize}[leftmargin=*,itemsep=1pt,topsep=2pt]
\item A formulation of numerical desiderata for embedding spaces, including what symmetric cosine geometry cannot
express (\S\ref{sec:form}).
\item A verified-reasoner distillation protocol with five objectives and three representation choices: a single space,
a learned numerical projection, and a parser hybrid (\S\ref{sec:method}).
\item An experimental design whose controls target C1--C3: hard-negative discrimination, a budget-matched reasoner-free
ablation, a parser-free stress slice and quantity-level leakage control (\S\ref{sec:design}).
\item Pre-declared hypotheses and decision rules with preliminary results (\S\ref{sec:results}), and a discussion of
where LLM reasoners can \emph{genuinely} add information (\S\ref{sec:disc}).
\end{itemize}

% =============================================================================================
\section{Related Work}
% =============================================================================================
\textbf{Numeracy in learned representations.} Naik et al.~\cite{naik2019numeracy} and Wallace et
al.~\cite{wallace2019numbers} probe word and contextual embeddings for magnitude and ordering. They find partial but
brittle numeracy that degrades outside the training range. Zhang et al.~\cite{zhang2020scales} show that embeddings
encode physical scales only coarsely, and NumerSense~\cite{lin2020numersense} probes numerical commonsense in masked
language models. Thawani et al.~\cite{thawani2021survey} survey number representations and argue for numeracy-aware
design.

\textbf{Numeracy-preserving representations.} Explicit numeric encodings such as DICE~\cite{sundararaman2020dice}
build numerical distance into the embedding by construction. Language-modelling work improves number
prediction~\cite{spithourakis2018numeracy}, and Geva et al.~\cite{geva2020injecting} inject numerical reasoning skills
through synthetic data. These approaches target token-level or generative representations. We target
\emph{sentence-level retrieval embeddings}, where numbers must coexist with context and meaning.

\textbf{Sentence embeddings and distillation.} Sentence-BERT~\cite{reimers2019sbert} and contrastive objectives such
as SimCSE~\cite{gao2021simcse} produce the encoders deployed today, and MiniLM~\cite{wang2020minilm} is itself a
distilled model. Reimers and Gurevych~\cite{reimers2020multilingual} distil an embedding space into a encoder while
preserving it, which is the template for our semantic-preservation term. Benchmarks such as
MTEB~\cite{muennighoff2023mteb} and STS~\cite{cer2017sts} measure general quality but contain few numerically
discriminative cases.

\textbf{Preserving pretrained behaviour.} Fine-tuning risks catastrophic forgetting. Parameter-space regularisers such
as EWC~\cite{kirkpatrick2017ewc} and L2-SP~\cite{li2018l2sp} constrain weights; we instead constrain \emph{outputs},
distilling from the frozen pretrained encoder on general text.

\textbf{Probing methodology and order.} Linear probes measure decodability, not use, and high probe accuracy can
reflect probe capacity~\cite{hewitt2019control}, so we report geometric and probe metrics separately. Order is
antisymmetric, and order embeddings~\cite{vendrov2016order} encode it with asymmetric scores; this matters because
cosine is symmetric (\S\ref{sec:form}).

\textbf{Positioning.} Prior work either \emph{probes} numeracy or \emph{adds} explicit numeric channels. We study
\emph{transfer} from an LLM reasoner into a retrieval encoder. The distinctive element is the set of controls that
attribute gains to the reasoner, the parser or the learned geometry.

% =============================================================================================
\section{Numerical Desiderata for Embedding Spaces}
\label{sec:form}
% =============================================================================================
Each statement $x$ expresses a quantity $\kappa(x)=(d_x,v_x)$: a dimension $d_x$ (mass, length, time, temperature,
volume, proportion) and a value $v_x>0$ in a canonical unit (g, m, s, K, L, ratio). Temperature is \emph{affine}
($v=T_{\mathrm C}+273.15$; $v=(T_{\mathrm F}-32)\tfrac59+273.15$), so ratios and log distances are always taken on the
absolute scale; Celsius or Fahrenheit ratios are never formed. With normalised embeddings $e(x)$ and
$s(x,y)=e(x)^\top e(y)$, we ask an embedding to satisfy five desiderata:
\begin{description}[leftmargin=0pt,itemsep=1pt,topsep=2pt]
\item[D1 Equivalence.] Equal quantities (10\,kg and 10{,}000\,g; 10 and 10.00, but \emph{not} 10.01) score above
non-equivalent pairs of the same dimension.
\item[D2 Graded distance.] For same-dimension pairs, $1-s$ grows monotonically with $|\log_{10}(v_x/v_y)|$. We also
study relative absolute distance $|v_x-v_y|/\max(v_x,v_y)$. The two disagree, and which is ``right'' depends on the
application.
\item[D3 Order.] The ordering $v_x<v_y$ is recoverable from the representation.
\item[D4 Dimensional separation.] Pairs of different dimensions (10\,kg, 10\,m) score low; no distance between them is
meaningful.
\item[D5 Context and semantics.] ``The patient's temperature is 37\,\textdegree C'' and ``the machine operates at
37\,\textdegree C'' are numerically equal but stay distinct, and general semantic similarity does not degrade.
\end{description}

\textbf{Three tensions.} (i) \emph{D3 cannot be expressed by cosine}: $s$ is symmetric, so order appears at best as
betweenness ($s(a,c)<\min(s(a,b),s(b,c))$ for $v_a<v_b<v_c$) or along a probed direction. (ii) \emph{D1 competes with
D5}: the same quantity in two contexts must be numerically equal yet semantically different. (iii) \emph{D2 competes
with D1 near the top of the scale}: pushing close pairs towards $s=1$ blurs exact equivalence, so graded targets are
capped at $c_{\max}<1$. These tensions suggest a single cosine space may be the wrong object; we therefore also test a
separate numerical subspace and a parser-hybrid similarity.

% =============================================================================================
\section{Method}
\label{sec:method}
% =============================================================================================
\textbf{Verified reasoner.} Qwen2.5-1.5B-Instruct~\cite{qwen2025qwen25} answers, per statement pair: same dimension
(Yes/No), equivalent after conversion (Yes/No), which is larger (A/B/Equal) and, for a random half of the same-dimension pairs, the ratio $A/B$ (a number); and, per
statement, whether the value is plausible in context (Yes/No). Categorical answers are read from the next-token
distribution over option tokens in one forward pass, giving a label and a confidence $p$; ratios use greedy decoding.
Every determinable answer is checked against Python ground truth; LLM output is never ground truth.

\textbf{Objectives.} The encoder all-MiniLM-L6-v2 minimises
$\Lq=\lambda_1\Lq_{\mathrm{eq}}+\lambda_2\Lq_{\mathrm{ord}}+\lambda_3\Lq_{\mathrm{dist}}+\lambda_4\Lq_{\mathrm{ctx}}+\lambda_5\Lq_{\mathrm{sem}}$.
\textbf{Equivalence} uses a confidence-weighted symmetric InfoNCE~\cite{oord2018cpc} over same-context equivalent pairs
$(a_i,b_i)$ with weights $w_i$ ($p$, or 1 for direct supervision):
\begin{equation}
\Lq_{\mathrm{eq}}=-\tfrac{1}{2\sum_i w_i}\textstyle\sum_i w_i\big[\log\sigma_j(s_{a_ib_j}/\tau)_i+\log\sigma_j(s_{b_ia_j}/\tau)_i\big],
\end{equation}
with $\sigma_j$ the softmax over in-batch partners. \textbf{Ordering} combines a margin ranking loss on
a learned direction $r(x)=u^\top e(x)$ with a betweenness loss on ordered triplets
$v_{\mathrm{lo}}<v_{\mathrm{mid}}<v_{\mathrm{hi}}$ (writing $s_{\mathrm{lh}}=s(x_{\mathrm{lo}},x_{\mathrm{hi}})$, etc.):
\begin{align}
\Lq_{\mathrm{ord}}&=\overline{[m_r-r(x_{\mathrm{hi}})+r(x_{\mathrm{lo}})]_+}\nonumber\\
&\quad+\overline{[s_{\mathrm{lh}}-s_{\mathrm{lm}}+m_t]_+ +[s_{\mathrm{lh}}-s_{\mathrm{mh}}+m_t]_+}.
\end{align}
\textbf{Distance} regresses cosine onto a magnitude-similarity target for same-dimension, non-equivalent pairs:
\begin{align}
\Lq_{\mathrm{dist}}&=\overline{\big(s(a,b)-t(\Delta)\big)^2},\\
t(\Delta)&=c_{\min}+(c_{\max}-c_{\min})\,e^{-\Delta/\sigma},\nonumber
\end{align}
with $\Delta=|\log_{10}(v_a/v_b)|$ (log mode) or $|v_a-v_b|/\max(v_a,v_b)$ (abs mode); under distillation, $\Delta$ is
computed from the reasoner's ratio. \textbf{Context} places upper bounds on similarity,
\begin{equation}
\Lq_{\mathrm{ctx}}=\overline{[s(a,b)-m_k]_+},
\end{equation}
where $m_k$ depends on whether the pair has different dimensions, is the same quantity in a different context, or is a
plausible versus contradictory statement. \textbf{Semantic preservation} distils from the frozen pretrained encoder
$f_0$ on unlabeled general sentences $g$~\cite{reimers2020multilingual}:
\begin{equation}
\Lq_{\mathrm{sem}}=\overline{1-e(g)^\top f_0(g)}.
\end{equation}

\textbf{Representations.}
\begin{itemize}[leftmargin=*,itemsep=1pt,topsep=2pt]
\item \emph{Single space}: all objectives act on $e$; similarity is $s$.
\item \emph{Learned numerical projection}: a linear head $P\!:\!384\!\to\!128$ receives the numerical objectives on
$z=\mathrm{norm}(Pe)$, while $\Lq_{\mathrm{sem}}$ acts on $e$. Similarity is $\alpha s+(1-\alpha)z^\top z'$. No parser is
involved.
\item \emph{Parser hybrid}: a regex parser extracts $(d,v)$ and gives $s_{\mathrm{num}}$: 1 if equivalent,
$c_{\max}e^{-\Delta/\sigma}$ if same dimension, 0 otherwise. Similarity is $\alpha s+(1-\alpha)s_{\mathrm{num}}$,
falling back to $s$ when parsing fails.
\end{itemize}
$\alpha$ is tuned on validation in both cases.

\textbf{Training.} AdamW (lr $2{\times}10^{-5}$, weight decay 0.01, 10\,\% warm-up), $4\times300$ steps, batch 32 pairs per
objective (16 triplets), $\lambda=(1,0.5,1,0.5,1)$, $\tau=0.05$, $m_r=0.1$, $m_t=0.05$, $m_k\in\{0.3,0.85,0.7\}$,
$(c_{\min},c_{\max})=(0.2,0.9)$, $\sigma_{\log}=1$, $\sigma_{\mathrm{abs}}=0.25$; checkpoint selection by the unweighted
validation loss, identical across runs; seeds 42--44; reasoner budget 4{,}000 training pairs.

% =============================================================================================
\section{Benchmark and Experimental Design}
\label{sec:design}
% =============================================================================================
\textbf{Data.} A deterministic generator (seed 42) produces \textbf{5{,}010 statements} in seven categories
(Table~\ref{tab:data}) over six dimensions, with decimals, negative temperatures, thousands separators, scientific
notation and metric/imperial units, abbreviated or spelled out.
\emph{Leakage control}: statements are keyed by quantity $(d,v)$; all statements of one quantity---any unit, precision,
context or category---share an equivalence group and a split, and multi-quantity sets are placed whole. Seven
templates appear only in test (76 statements). The split is 3{,}541/761/708, with zero groups crossing splits; pairs
and triplets are built within splits. The test set has 4{,}135 pairs (351 equivalent same-context, 327 same quantity
in a different context, 1{,}134 close $\le2\times$, 597 mid, 1{,}047 distant $\ge10\times$, 642 cross-dimension, 37
plausibility) and 2{,}600 ordering triplets.
\emph{Parser-free slice}: 379 test statements are rewritten with spelled-out integers (``thirty-five degrees
Celsius''); the parser reads none, so gains there must come from the embedding.

\begin{table}[t]
\centering\footnotesize
\caption{Benchmark categories (ground truth computed in Python).}
\label{tab:data}
\begin{tabular}{@{}clr@{}}
\toprule
Cat. & Content & $n$ \\
\midrule
A & Magnitude (\textit{The cable measures 0.75 miles.}) & 900 \\
B & Unit-equivalence sets & 1{,}202 \\
C & Ordering triplets & 601 \\
D & Same quantity, several contexts & 602 \\
E & Precision (10 / 10.0 / 10.00 / 10.01 / 10.1) & 500 \\
F & Ratios, percentages, differences & 604 \\
G & Plausible vs contradictory (\textit{0.006 mL syrup dose}) & 601 \\
\bottomrule
\end{tabular}
\end{table}

\textbf{Configurations} are listed in Table~\ref{tab:configs}. Three are controls for the confounds:
\vMatched trains on \emph{exactly} the reasoner's pairs with ground-truth labels, so \vVerified$-$\vMatched isolates what
the reasoner adds beyond pair selection and confidence weights (C3); \vDual learns a numerical subspace without
any parser; and the \textbf{words slice} disables the parser at test time (C2). To keep this comparison fair, \vMatched mirrors the reasoner's supervision exactly: it receives distance targets only on the pairs for which the reasoner was asked a ratio (half of the same-dimension pairs), and neither it nor the reasoner-based models use the triplet betweenness term, which only \vExact receives. The same-/different-context labels used to form positives come from the benchmark metadata for every model, not from the reasoner.

\begin{table}[t]
\centering\footnotesize
\caption{Configurations (identical test data for all).}
\label{tab:configs}
\begin{tabular}{@{}l>{\raggedright\arraybackslash}p{4.4cm}@{}}
\toprule
\vBase & Original all-MiniLM-L6-v2, cosine. \\
\vParse & Original embeddings + parser hybrid. \\
\vParseOnly & Parser only (upper reference on digit text). \\
\vExact & Direct mathematical supervision, all training pairs. \\
\vMatched & Direct supervision on the reasoner's pair budget (\textbf{reasoner-free control}). \\
\vRaw & Same pairs, raw reasoner labels (confidence-weighted). \\
\vVerified & Same pairs, reasoner labels kept only where verified. \\
\vExactParse & \vExact embeddings + parser hybrid. \\
\vDual & Learned numerical projection head, no parser. \\
\vAbs & \vExact with the absolute-distance target (seed 42 only). \\
Ablation & $\Lq_{\mathrm{sem}}$ only, then $+\Lq_{\mathrm{eq}}$, $+\Lq_{\mathrm{ord}}$, $+\Lq_{\mathrm{dist}}$, $+\Lq_{\mathrm{ctx}}$ (= \vExact). \\
\bottomrule
\end{tabular}
\end{table}

\textbf{Metrics.} \emph{Equivalence}: ROC-AUC / balanced accuracy of equivalent vs \emph{hard} negatives (same
dimension, $\le10\times$ apart; C1), with cross-unit and precision subsets. \emph{Order}: per-dimension ridge probe to
$\log_{10}v$ (pairwise accuracy) and betweenness under each model's similarity $s$. \emph{Distance}: Spearman $\rho$ of $1-s$ vs $\Delta$ on
same-template pairs. \emph{Retrieval}: MRR for same-group retrieval among all test statements. \emph{Context}:
dimension discrimination $\Pr[s(a,\text{equiv.})>s(a,\text{cross-dim.})]$ and linear context and plausibility probes.
\emph{Semantics}: STS-B~\cite{cer2017sts} Spearman. Thresholds and $\alpha$ are tuned on validation.

\textbf{Decision rule and hypotheses.} A difference is \emph{supported} if it exceeds $\max(0.01,2\hat\sigma)$
($\hat\sigma$: seed standard deviation), \emph{contradicted} if it is a decrease of that size, and \emph{inconclusive}
otherwise. We declare, before seeing results:
\begin{description}[leftmargin=0pt,itemsep=1pt,topsep=2pt]
\item[H1] The reasoner is reliable on dimension identification but less reliable on cross-unit equivalence and ratios.
\item[H2] The original encoder fails D1 and D2: its cosine is high for all pair types, so hard-negative equivalence AUC
is near chance and distance correlation is weak.
\item[H3] Direct supervision (\vExact) improves D1, D2 and D4 over \vBase on digit-form text.
\item[H4] Verified distillation (\vVerified) does not exceed budget-matched direct supervision (\vMatched), because of the
verification paradox; raw labels (\vRaw) do worse.
\item[H5] Parser hybrids (\vParse, \vExactParse) dominate on digit-form text but gain nothing on the words slice.
\item[H6] Semantic preservation keeps the STS-B drop within 0.02.
\end{description}

% =============================================================================================
\section{Preliminary Results}
\label{sec:results}
% =============================================================================================
Hardware: \envv{hardware}; reasoner \teach{teacher_model}; means $\pm$ seed standard deviation where available.

\textbf{Reasoner reliability (H1).} Table~\ref{tab:reasoner} reports the reasoner against Python ground truth.
The reasoner is far from reliable. Its equivalence judgements carry some signal (soft AUC \teach{equivalence_auc}),
but raw accuracy (\teach{equivalence_acc}) barely exceeds the majority baseline (\teach{majority_baseline_equivalence}) and
balanced accuracy is only \teach{equivalence_balanced_acc}. It recognises only \teach{unit_conversion_acc_cross_unit_equiv} of
cross-unit equivalences and places only \teach{ratio_acc_within_1pct} of ratios within 1\,\%, although every ratio answer
parses as a number. Its dimension and order answers collapse onto a single option: it answers ``Yes'' to the
same-dimension question for only 3\,\% of pairs, although 84\,\% share a dimension (accuracy \teach{same_dimension_acc}), and it
places 89\,\% of its order probability on ``Equal'' (accuracy \teach{ordering_acc}, close to the share of equal pairs).
Only \teach{fully_verified_fraction} of the pairs survive verification. H1 is therefore only \emph{partly} supported: the
reasoner is weak on conversion and ratios, as predicted, but it is \emph{not} reliable on dimension identification.
\textbf{This is a property of the model, not of the scoring.} A diagnostic confirms that 95--100\,\% of the
next-token probability falls on the option tokens, that free-text answers agree with the scored ones (``No'' for
10\,kg vs 10{,}000\,g; ``Equal'' for 2\,m vs 300\,cm), and that batching changes scores only at the third decimal.
Rewording the questions on 150 training pairs shifts \emph{which} option the model prefers (order answers move from
``Equal'' to mostly ``B'') but not its accuracy: 0.18 to 0.22 on dimension (majority baseline 0.86) and 0.31 to 0.33
on order (majority 0.39). A 1.5B reasoner therefore supplies little usable supervision on determinable relations,
which bears directly on H4.

\begin{table}[t]
\centering\footnotesize
\caption{Reasoner accuracy vs ground truth (\teach{n_pairs} pairs).}
\label{tab:reasoner}
\begin{tabular}{@{}lr@{\hspace{1.2em}}lr@{}}
\toprule
Same dimension & \teach{same_dimension_acc} & Ordering & \teach{ordering_acc} \\
Equiv.\ (balanced) & \teach{equivalence_balanced_acc} & Ratio within 1\,\% & \teach{ratio_acc_within_1pct} \\
Equiv.\ AUC (soft) & \teach{equivalence_auc} & Plausibility (bal.) & \teach{compatibility_balanced_acc} \\
Cross-unit equiv.\ (recall) & \teach{unit_conversion_acc_cross_unit_equiv} & Fully verified & \teach{fully_verified_fraction} \\
\bottomrule
\end{tabular}
\end{table}

\textbf{What the original encoder sees (H2).} Table~\ref{tab:cos} gives the mean cosine by pair type.
\interp{is the \vBase column flat across pair types? If so, C1 holds.}

\begin{table}[t]
\centering\footnotesize
\caption{Mean cosine by test pair type (\vExact: seed 42).}
\label{tab:cos}
\begin{tabular}{@{}lrcc@{}}
\toprule
Pair type & $n$ & \vBase & \vExact \\
\midrule
Equivalent, same context & \res{\Ea}{cosn.equivalent} & \res{\Ea}{cos.equivalent} & \res{\Ec}{cos.equivalent} \\
Same qty, other context & \res{\Ea}{cosn.equivalent_diff_context} & \res{\Ea}{cos.equivalent_diff_context} & \res{\Ec}{cos.equivalent_diff_context} \\
Close ($\le2\times$) & \res{\Ea}{cosn.close} & \res{\Ea}{cos.close} & \res{\Ec}{cos.close} \\
Distant ($\ge10\times$) & \res{\Ea}{cosn.distant} & \res{\Ea}{cos.distant} & \res{\Ec}{cos.distant} \\
Different dimension & \res{\Ea}{cosn.cross_dimension} & \res{\Ea}{cos.cross_dimension} & \res{\Ec}{cos.cross_dimension} \\
Plausible vs contradictory & \res{\Ea}{cosn.plausibility} & \res{\Ea}{cos.plausibility} & \res{\Ec}{cos.plausibility} \\
\bottomrule
\end{tabular}
\end{table}

\begin{table*}[t]
\centering\footnotesize
\caption{Numerical metrics on digit-form test text. AUC: equivalent vs hard negatives; Order: probe pairwise
accuracy; Betw.: betweenness under the model's similarity; $\rho_{\log}$: Spearman of $1-s$ vs $|\Delta\log_{10}|$ (same template).}
\label{tab:num}
\setlength{\tabcolsep}{4pt}\begin{tabular}{@{}lccccccc@{}}
\toprule
Model & Equiv.\ AUC & Unit conv. & Precision & Order (probe) & Betw. & $\rho_{\log}$ & MRR \\
\midrule
\vBase & \respm{\Ea}{equivalence.equivalence_auc_hard} & \respm{\Ea}{equivalence.unit_conversion_bal_acc} & \respm{\Ea}{equivalence.precision_bal_acc} & \respm{\Ea}{ordering.probe_pairwise_order_acc} & \respm{\Ea}{ordering.triplet_betweenness_acc} & \respm{\Ea}{distance.same_template_log_distance_spearman} & \respm{\Ea}{retrieval.mrr} \\
\vParse & \respm{\Eb}{equivalence.equivalence_auc_hard} & \respm{\Eb}{equivalence.unit_conversion_bal_acc} & \respm{\Eb}{equivalence.precision_bal_acc} & \respm{\Eb}{ordering.probe_pairwise_order_acc} & \respm{\Eb}{ordering.triplet_betweenness_acc} & \respm{\Eb}{distance.same_template_log_distance_spearman} & \respm{\Eb}{retrieval.mrr} \\
\vParseOnly & \respm{\Ep}{equivalence.equivalence_auc_hard} & \respm{\Ep}{equivalence.unit_conversion_bal_acc} & \respm{\Ep}{equivalence.precision_bal_acc} & \respm{\Ep}{ordering.probe_pairwise_order_acc} & \respm{\Ep}{ordering.triplet_betweenness_acc} & \respm{\Ep}{distance.same_template_log_distance_spearman} & \respm{\Ep}{retrieval.mrr} \\
\midrule
\vExact & \respm{\Ec}{equivalence.equivalence_auc_hard} & \respm{\Ec}{equivalence.unit_conversion_bal_acc} & \respm{\Ec}{equivalence.precision_bal_acc} & \respm{\Ec}{ordering.probe_pairwise_order_acc} & \respm{\Ec}{ordering.triplet_betweenness_acc} & \respm{\Ec}{distance.same_template_log_distance_spearman} & \respm{\Ec}{retrieval.mrr} \\
\vMatched & \respm{\Ecb}{equivalence.equivalence_auc_hard} & \respm{\Ecb}{equivalence.unit_conversion_bal_acc} & \respm{\Ecb}{equivalence.precision_bal_acc} & \respm{\Ecb}{ordering.probe_pairwise_order_acc} & \respm{\Ecb}{ordering.triplet_betweenness_acc} & \respm{\Ecb}{distance.same_template_log_distance_spearman} & \respm{\Ecb}{retrieval.mrr} \\
\vRaw & \respm{\Eda}{equivalence.equivalence_auc_hard} & \respm{\Eda}{equivalence.unit_conversion_bal_acc} & \respm{\Eda}{equivalence.precision_bal_acc} & \respm{\Eda}{ordering.probe_pairwise_order_acc} & \respm{\Eda}{ordering.triplet_betweenness_acc} & \respm{\Eda}{distance.same_template_log_distance_spearman} & \respm{\Eda}{retrieval.mrr} \\
\vVerified & \respm{\Edb}{equivalence.equivalence_auc_hard} & \respm{\Edb}{equivalence.unit_conversion_bal_acc} & \respm{\Edb}{equivalence.precision_bal_acc} & \respm{\Edb}{ordering.probe_pairwise_order_acc} & \respm{\Edb}{ordering.triplet_betweenness_acc} & \respm{\Edb}{distance.same_template_log_distance_spearman} & \respm{\Edb}{retrieval.mrr} \\
\vExactParse & \respm{\Ee}{equivalence.equivalence_auc_hard} & \respm{\Ee}{equivalence.unit_conversion_bal_acc} & \respm{\Ee}{equivalence.precision_bal_acc} & \respm{\Ee}{ordering.probe_pairwise_order_acc} & \respm{\Ee}{ordering.triplet_betweenness_acc} & \respm{\Ee}{distance.same_template_log_distance_spearman} & \respm{\Ee}{retrieval.mrr} \\
\vDual & \respm{\Ef}{equivalence.equivalence_auc_hard} & \respm{\Ef}{equivalence.unit_conversion_bal_acc} & \respm{\Ef}{equivalence.precision_bal_acc} & \respm{\Ef}{ordering.probe_pairwise_order_acc} & \respm{\Ef}{ordering.triplet_betweenness_acc} & \respm{\Ef}{distance.same_template_log_distance_spearman} & \respm{\Ef}{retrieval.mrr} \\
\vAbs & \respm{\Eabs}{equivalence.equivalence_auc_hard} & \respm{\Eabs}{equivalence.unit_conversion_bal_acc} & \respm{\Eabs}{equivalence.precision_bal_acc} & \respm{\Eabs}{ordering.probe_pairwise_order_acc} & \respm{\Eabs}{ordering.triplet_betweenness_acc} & \respm{\Eabs}{distance.same_template_log_distance_spearman} & \respm{\Eabs}{retrieval.mrr} \\
\bottomrule
\end{tabular}
\end{table*}

\begin{table*}[t]
\centering\footnotesize
\caption{Context, semantics and the parser-free slice (numbers in words; hybrids fall back to their embedding).}
\label{tab:ctxwords}
\setlength{\tabcolsep}{4pt}\begin{tabular}{@{}lccccccc@{}}
\toprule
& \multicolumn{4}{c}{Context and semantics} & \multicolumn{3}{c}{Words slice vs digits} \\
\cmidrule(lr){2-5}\cmidrule(l){6-8}
Model & Dim.\ discrim. & Context probe & Plausibility probe & STS-B $\rho$ & MRR (digits) & MRR (words) & Equiv.\ AUC (words) \\
\midrule
\vBase & \respm{\Ea}{context.dimension_discrimination_acc} & \respm{\Ea}{context.context_probe_acc} & \respm{\Ea}{context.plausibility_probe_acc} & \respm{\Ea}{semantic.stsb_spearman} & \respm{\Ea}{retrieval.mrr} & \respm{\Ea}{stress_words.words_mrr} & \respm{\Ea}{stress_words.words_equivalence_auc} \\
\vParse & \respm{\Eb}{context.dimension_discrimination_acc} & \respm{\Eb}{context.context_probe_acc} & \respm{\Eb}{context.plausibility_probe_acc} & \respm{\Eb}{semantic.stsb_spearman} & \respm{\Eb}{retrieval.mrr} & \respm{\Eb}{stress_words.words_mrr} & \respm{\Eb}{stress_words.words_equivalence_auc} \\
\vExact & \respm{\Ec}{context.dimension_discrimination_acc} & \respm{\Ec}{context.context_probe_acc} & \respm{\Ec}{context.plausibility_probe_acc} & \respm{\Ec}{semantic.stsb_spearman} & \respm{\Ec}{retrieval.mrr} & \respm{\Ec}{stress_words.words_mrr} & \respm{\Ec}{stress_words.words_equivalence_auc} \\
\vVerified & \respm{\Edb}{context.dimension_discrimination_acc} & \respm{\Edb}{context.context_probe_acc} & \respm{\Edb}{context.plausibility_probe_acc} & \respm{\Edb}{semantic.stsb_spearman} & \respm{\Edb}{retrieval.mrr} & \respm{\Edb}{stress_words.words_mrr} & \respm{\Edb}{stress_words.words_equivalence_auc} \\
\vExactParse & \respm{\Ee}{context.dimension_discrimination_acc} & \respm{\Ee}{context.context_probe_acc} & \respm{\Ee}{context.plausibility_probe_acc} & \respm{\Ee}{semantic.stsb_spearman} & \respm{\Ee}{retrieval.mrr} & \respm{\Ee}{stress_words.words_mrr} & \respm{\Ee}{stress_words.words_equivalence_auc} \\
\vDual & \respm{\Ef}{context.dimension_discrimination_acc} & \respm{\Ef}{context.context_probe_acc} & \respm{\Ef}{context.plausibility_probe_acc} & \respm{\Ef}{semantic.stsb_spearman} & \respm{\Ef}{retrieval.mrr} & \respm{\Ef}{stress_words.words_mrr} & \respm{\Ef}{stress_words.words_equivalence_auc} \\
\bottomrule
\end{tabular}
\end{table*}

\textbf{Equivalence, order and distance (H3, H4).} Table~\ref{tab:num} is the main comparison and Fig.~\ref{fig:cos}
shows the similarity distributions. \interp{H3: \vExact vs \vBase per column. H4: \vVerified vs \vMatched and \vRaw vs \vMatched. Does
order move in the probe but not in betweenness, as \S\ref{sec:form} predicts?}

\begin{figure}[t]
\centering\figorbox{01_cosine_baseline_vs_finetuned.png}{\columnwidth}
\caption{Cosine similarity by numerical relation: original, directly fine-tuned and reasoner-distilled encoders.}
\label{fig:cos}
\end{figure}

\textbf{Context, semantics and attribution (H5, H6).} Table~\ref{tab:ctxwords} pairs context and STS-B metrics with
the parser-free slice. \interp{STS-B change (H6). Digit-to-words MRR drop of \vParse and \vExactParse vs \vExact and \vDual: a large hybrid
drop with stable \vExact and \vDual is evidence of learned numeracy; equal drops mean parser-driven gains (H5).}

\textbf{Retrieval and geometry.} Fig.~\ref{fig:retr} shows equivalence retrieval and the
embedding geometry, and Fig.~\ref{fig:dist} plots embedding against numerical distance. \interp{does magnitude become
a visible direction in the PCA after fine-tuning?}

\textbf{Ablation.} Table~\ref{tab:abl} adds the objectives cumulatively (seed 42). \interp{largest step; steps smaller
than the \vExact seed standard deviation are inconclusive.}

\begin{table}[t]
\centering\footnotesize\setlength{\tabcolsep}{3.5pt}
\caption{Cumulative ablation (seed 42); the last row is \vExact.}
\label{tab:abl}
\begin{tabular}{@{}lcccccc@{}}
\toprule
Variant & AUC & MRR & $\rho_{\log}$ & Order & Dim. & STS-B \\
\midrule
$\Lq_{\mathrm{sem}}$ only & \res{abl_A_semantic_only}{equivalence.equivalence_auc_hard} & \res{abl_A_semantic_only}{retrieval.mrr} & \res{abl_A_semantic_only}{distance.same_template_log_distance_spearman} & \res{abl_A_semantic_only}{ordering.probe_pairwise_order_acc} & \res{abl_A_semantic_only}{context.dimension_discrimination_acc} & \res{abl_A_semantic_only}{semantic.stsb_spearman} \\
$+\,\Lq_{\mathrm{eq}}$ & \res{abl_B_+equivalence}{equivalence.equivalence_auc_hard} & \res{abl_B_+equivalence}{retrieval.mrr} & \res{abl_B_+equivalence}{distance.same_template_log_distance_spearman} & \res{abl_B_+equivalence}{ordering.probe_pairwise_order_acc} & \res{abl_B_+equivalence}{context.dimension_discrimination_acc} & \res{abl_B_+equivalence}{semantic.stsb_spearman} \\
$+\,\Lq_{\mathrm{ord}}$ & \res{abl_C_+ordering}{equivalence.equivalence_auc_hard} & \res{abl_C_+ordering}{retrieval.mrr} & \res{abl_C_+ordering}{distance.same_template_log_distance_spearman} & \res{abl_C_+ordering}{ordering.probe_pairwise_order_acc} & \res{abl_C_+ordering}{context.dimension_discrimination_acc} & \res{abl_C_+ordering}{semantic.stsb_spearman} \\
$+\,\Lq_{\mathrm{dist}}$ & \res{abl_D_+distance}{equivalence.equivalence_auc_hard} & \res{abl_D_+distance}{retrieval.mrr} & \res{abl_D_+distance}{distance.same_template_log_distance_spearman} & \res{abl_D_+distance}{ordering.probe_pairwise_order_acc} & \res{abl_D_+distance}{context.dimension_discrimination_acc} & \res{abl_D_+distance}{semantic.stsb_spearman} \\
$+\,\Lq_{\mathrm{ctx}}$ (full) & \res{\Ec}{equivalence.equivalence_auc_hard} & \res{\Ec}{retrieval.mrr} & \res{\Ec}{distance.same_template_log_distance_spearman} & \res{\Ec}{ordering.probe_pairwise_order_acc} & \res{\Ec}{context.dimension_discrimination_acc} & \res{\Ec}{semantic.stsb_spearman} \\
\bottomrule
\end{tabular}
\end{table}

\begin{figure*}[t]
\centering
\begin{minipage}{0.49\textwidth}\centering\figorbox{03_equivalence_retrieval.png}{\linewidth}\end{minipage}\hfill
\begin{minipage}{0.49\textwidth}\centering\figorbox{06_pca_embedding_space.png}{\linewidth}\end{minipage}
\caption{Left: equivalence retrieval (R@$k$, MRR, MRR on the words slice). Right: PCA of mass statements coloured by
$\log_{10}$ magnitude, original vs fine-tuned encoder.}
\label{fig:retr}
\end{figure*}

\begin{figure}[t]
\centering\figorbox{02_numerical_vs_embedding_distance.png}{\columnwidth}
\caption{Numerical log distance vs embedding distance (same-template pairs).}
\label{fig:dist}
\end{figure}

% =============================================================================================
\section{Discussion}
\label{sec:disc}
% =============================================================================================
\textbf{Higher cosine is not better numeracy.} A general encoder can give high similarity to every pair of
template-sharing sentences, including pairs four orders of magnitude apart or in different dimensions. Reporting mean
cosine on positives alone mistakes this compression for understanding. The informative quantities are discriminative
(AUC against hard negatives), graded (rank correlation with numerical distance) and comparative (equivalent vs
cross-dimension partners of one anchor). Fine-tuning may \emph{lower} all absolute cosines while improving each of
these---the signature of a space that has begun to separate numbers, not of degradation.

\textbf{The verification paradox.} Verification is necessary: an unverified small reasoner injects noise into exactly
the relations it should teach. But it also bounds the reasoner's value on determinable relations, where filtered labels
coincide with a few lines of Python. What remains is a \emph{selection effect} (the pairs the reasoner gets right may be
easier) and \emph{confidence weights}. Without the budget-matched \vMatched, any \vVerified gain over \vBase would be wrongly
credited to distillation. If H4 holds, the honest conclusion is that \emph{LLM reasoners are the wrong source of
supervision for relations a parser can compute}. Their value lies in signals Python cannot verify: plausibility (is
37\,\textdegree C a normal body temperature?), implicit units (``a six-footer''), approximations (``about a dozen''),
ranges and free phrasing.

\textbf{Learned or parsed?} A parser hybrid will look near-perfect whenever numbers appear as digits with canonical
units---a property of the benchmark, not evidence of learning. The words slice breaks this. If the hybrids' advantage
vanishes there while a fine-tuned or projected embedding keeps part of its gain, the embedding has internalised some
numeracy. If everything collapses to baseline, the encoder learned digit-string regularities rather than quantities,
pointing to number-aware tokenisation~\cite{thawani2021survey,sundararaman2020dice} as the bottleneck.

\textbf{One space or two; log or absolute distance?} The tensions of \S\ref{sec:form} argue for decoupling: a
numerical subspace (\vDual) can collapse equivalent quantities while the semantic space keeps contexts apart, and order
may ultimately need an asymmetric score~\cite{vendrov2016order}. Log distance suits magnitude reasoning; relative
absolute distance suits tolerance-based matching in entity resolution. The choice should follow the downstream
decision, and benchmarks should report both.

\textbf{Semantic cost.} Self-distillation on general text~\cite{reimers2020multilingual} is a cheap guard
against forgetting, but the numerical objectives deliberately reshape similarities between near-identical sentences,
so some movement on STS-B is expected. The relevant question is the size of the trade-off, not whether one exists.

\section{Limitations and Threats to Validity}
\textbf{Data.} The data are synthetic and template-based; real text has implicit units, ranges and several
quantities per sentence. Outside category G, values are deliberately decoupled from plausibility. \textbf{External validity.} Templates are shared across splits by design, except for seven held-out ones, and the
words slice tests format shift rather than new vocabulary. Conclusions about a 1.5B reasoner need not transfer to larger
models. \textbf{Internal validity.} Hyperparameters ($\lambda$, margins, $\sigma$) were fixed a priori rather than tuned
per configuration, so a different weighting could change the ablation ranking. \textbf{Construct validity.} Probes
measure linear decodability, not downstream use~\cite{hewitt2019control}. \textbf{Statistics.} Three seeds for the main
configurations and one for the ablation; bootstrap intervals cover test-item sampling, not training variance.

% =============================================================================================
\section{Conclusion}
% =============================================================================================
Transferring numerical reasoning into sentence embeddings is attractive and easy to evaluate badly. We argued that
claims about it require discriminative rather than absolute similarity, a parser-free test slice and a budget-matched
reasoner-free baseline, and we provided a benchmark and protocol implementing them with pre-declared hypotheses.
\interp{summary of which hypotheses were supported.} Promising directions are LLM reasoners for \emph{non-determinable}
numerical judgements, asymmetric heads for order, number-aware tokenisation, and evaluation on real numeric retrieval
and entity-resolution tasks. Code and configuration:
\url{https://github.com/TariqMahmood93/Numerical_Intelligence_Transfor}.

\bibliographystyle{IEEEtran}
\bibliography{refs}

\appendix[Reasoner Prompts]
System: \emph{``You are a careful assistant for physical quantities. Convert units exactly before comparing.
Temperatures must be compared on an absolute scale. Answer with only the requested option.''} User prompts present
``Statement A: \ldots\ Statement B: \ldots'' and ask:
\textbf{dimension}---``Do A and B measure the same physical dimension? Answer Yes or No.'';
\textbf{equivalence}---``After converting to a common unit, do A and B state exactly the same quantity? Answer Yes or
No.'';
\textbf{order}---``After converting to a common unit, which quantity is larger? Answer A, B, or Equal.'';
\textbf{ratio}---``Convert both quantities to the same unit (use kelvin for temperatures) and compute the ratio A
divided by B. Reply with only the number.'';
\textbf{plausibility}---``Is the numerical value physically plausible given the description in the statement? Answer
Yes or No.''

\end{document}
